\section{Introduction}\label{sec:intro}
Many distributed systems require two network endpoints holding sets $A$ and $B$ to find their symmetric difference $A\triangle B$, the elements present on only one side. This problem is called set reconciliation. The Invertible Bloom Lookup Table (IBLT)\cite{gm,eppstein2011whatsdifference} is the standard data structure for this task. Each side builds a table of $M$ cells and hashes each element into the count, keySum, and checksum fields of $k$ cells. The receiver then subtracts the tables cell by cell. Shared elements contribute identical values to the same cells and cancel exactly, leaving only the contributions of the $d=|A\triangle B|$ differing elements. An element present only in $A$ has direction $+1$, and an element present only in $B$ has direction $-1$. When the difference size $d$ is small enough relative to the table size $M$ (the load $d/M$ is below the peeling threshold), the receiver repeatedly finds a pure cell containing a single differing element and recovers all differing elements through peeling. Once $d$ exceeds the table's capacity, peeling stalls and decoding fails.

An IBLT presents a basic tension: the table size $M$ must be fixed at construction, yet the difference size $d$ remains unknown until synchronization completes. In practice, $d$ for the same pair of nodes can vary by several orders of magnitude over time. During routine synchronization, the two sides are nearly identical and $d$ is small. After recovery from a network partition or a batch write, $d$ can grow sharply. Allocating a large table for the worst case wastes bandwidth in most routine synchronizations. Allocating a small table causes decoding failure when the difference is large. The two-round protocol\cite{wu2026selfsizing} addresses this tension. The first round sends a small table of fixed size, allowing routine small differences to complete in one round. If an oversized difference causes decoding failure, the count array of the difference table still retains the superimposed contributions of all $d$ differing elements. The receiver estimates an upper bound $U$ on $d$ from the sum of squared counts in this table. The second round allocates sufficient capacity based on $U$ and completes synchronization. The key requirement is to control the underestimate risk conditioned on failure (the conditional miss rate). If $U<d$ (a miss), insufficient capacity puts the second-round table at risk of another failure. An overly conservative $U$ wastes second-round bandwidth.

The method in \cite{wu2026selfsizing} uses the count sum and the sum of squared counts in the first-round difference table to construct an unbiased estimator $\widehat d$ and obtain an upper bound $U$ by inversion. Estimating $d$ after failure is difficult because decoding failure itself selects the sample, and this selection couples with the sign composition of the set difference (the ratio of differing elements from $A$ to those from $B$). When we estimate $d$ with this sum-of-squares method, each element's own contribution is fixed. The estimation error comes entirely from cross terms between differing elements that fall into the same cell. The standard IBLT uses each element's set-membership direction ($+1$ or $-1$) as its accumulation sign in count. Two elements with the same direction add when they share a cell, increasing the sum of squares. Elements with opposite directions cancel. Peeling fails because the position mapping of differing elements to cells has too much overlap. Together, these effects produce systematic bias in the failure sample. For a one-sided difference, all differing elements come from the same side. This is common in one-way synchronization, for example when a replica lacks only data newly written to the primary. All cross terms then have the same sign, so the high overlap selected by failure is amplified in the sum of squared counts, substantially biasing the estimate. In two-way synchronization where each side writes new data, the numbers of differing elements on the two sides are comparable. For this balanced difference, positive and negative cross terms largely cancel, leaving the post-failure mean estimate nearly unchanged. The bias magnitude varies continuously with the ratio of differing elements on the two sides, and its sign reverses with the number of hashes $k$ (Section~\ref{sec:model}). After decoding failure, the receiver does not know whether the difference comes from one side or both. The observed sum of squared counts therefore cannot reveal how much is due to the true $d$ and how much to the direction bias induced by failure selection. Thus, except in deep overload, the upper-bound rule for the standard IBLT in \cite{wu2026selfsizing} requires offline simulation calibration for each configuration.

Our core insight is that separating an element's set direction from its accumulation sign removes the source of this coupling. We therefore propose the random-sign IBLT. The change is lightweight: the two sides share a random-sign hash independent of the position mapping, assigning each key a random $+1$ or $-1$. On insertion, the element's random sign multiplies its contributions to all three fields: count, keySum, and checksum.

This modification supports both decoding and estimation. Shared elements receive the same random sign in both tables and still cancel exactly after subtraction. A cell containing one unrecovered differing element has count $+1$ or $-1$. Multiplying keySum by count cancels the random sign exactly (since $(\pm1)^2=1$) and recovers the original key. Computing the key's random sign and comparing it with count then reveals whether the key belongs to set $A$ or $B$. Under ideal purity, the random-sign IBLT therefore peels exactly the same elements as the standard IBLT at every step. The decoding threshold, failure probability, and number and byte width of cell fields remain unchanged. For estimation, each differing element's coding coefficient in the difference table becomes an independent uniform random sign, whether the difference is one-sided or balanced. Peeling success depends only on which cells the elements select (the mapping), while cancellation of the estimation error depends on the random signs. These are independent. For every mapping selected by failure, the random signs remain independent and uniform, so their contributions cancel in expectation. Using random signs to cancel cross terms is a classical idea from the AMS sketch and CountSketch\cite{alon1999spacecomplexity,charikar2004countsketch}. We implement it in a peelable data structure, enabling one table to support both difference recovery and difference-size estimation and allowing rigorous statistical analysis of the post-failure estimator.

\paragraph{Main results and contributions.}
We prove that the random-sign IBLT preserves the standard IBLT's decoding behavior under ideal purity and that the joint distribution of the difference-table statistic and the mapping is identical for all sign compositions. Section~\ref{sec:design} presents the encoding and peeling algorithm that achieves this decoupling. We then prove that the estimator $\widehat d$ is exactly unbiased on every fixed mapping, giving conditional mean $d$ after failure, and fully characterize the conditional distribution of the estimation error across all load scales (Sections~\ref{sec:measurement} and \ref{sec:inversion}). For large differences, we derive a Gaussian rule from the conditional Gaussian limit. For small differences, we characterize a discrete limit distribution and construct a Skellam rule that uniformly controls the conditional miss rate for all difference sizes without offline simulation calibration. We also provide an efficient planted sampling method for rare failures at low load. The experiments (Section~\ref{sec:exp}) verify that random-sign and standard encoding have identical decoding behavior. They show that the Skellam rule keeps the conditional miss rate strictly within the target risk across the full load range spanning several orders of magnitude. With this rule, the two-round protocol reliably completes synchronization at a total transmission cost close to the ideal lower bound for known $d$.

\paragraph{Structure of the paper.}
Section~\ref{sec:model} introduces the standard IBLT model and analyzes its post-failure estimation bias. Section~\ref{sec:design} presents the random-sign IBLT and the two-round protocol. Section~\ref{sec:measurement} proves exact unbiasedness given the mapping and derives the Gaussian rule for large differences. Section~\ref{sec:inversion} derives the discrete limit distribution for small differences and the Skellam rule for all scales. Section~\ref{sec:exp} reports the experimental evaluation. Sections~\ref{sec:related} and \ref{sec:conclusion} discuss related work, rule selection, and computational cost. The appendix contains complete proofs and supplementary experiments.
